\subsection{節点型と $E_6$ 一般化DAHAの球面部分}
\label{subsec:cubic-CY-three-alpha-nodal-E6}

以下では節点型の命題\ref{prop:cubic-CY-three-alpha-nodal-secant-center}の記号を使い、
$q\ne0,\pm1$、$q^3\ne1$ とする。
$E_6$ 型の一般化DAHAを、Etingof--Oblomkov--Rains
\cite[第8.1節]{EtingofOblomkovRainsGDaha}の二生成元表示で書く。
同論文の変形係数の三乗をここでは $q$ とする。
従ってこの表示は
\[
 H=\frac{k\langle b,c\rangle}
 {(b^3-f_1b^2+f_2b-1,\ c^3-g_1c^2+g_2c-1, (bc)^3-e_2(bc)^2+e_1bc-q)}
\]
である。$c$ の定義多項式に単純根 $\lambda$ を一つ選び、
$g_2=g_1\lambda-\lambda^2+\lambda^{-1}$ とする。
その微分値を $\Delta=3\lambda^2-2g_1\lambda+g_2$、
対応するスペクトル冪等元を
\[
 e=\frac{c^2+(\lambda-g_1)c+\lambda^{-1}}{\Delta}
\]
と書く。

\begin{theorem}\label{thm:cubic-CY-three-alpha-nodal-E6}
この全単純根チャートで $eHe$ は節点型中央三次切断のアフィン環である。
係数写像は有理式で明示でき、その分母は
$\lambda q(q^2-1)(q^3-1)\Delta$ の因子に限られる。
$x,y,z$ の非零倍率を加えた係数写像は、各固定許容 $q$ で
$P_xP_yA_z\ne0$ の節点切断係数空間へ支配的である。
\end{theorem}
\begin{proof}
まず $X=ebe$、$Y=eb^{-1}e$、$Z=eb^{-1}cbe$ と置く。
$Z$ を $Z-f_2\lambda X-e_2Y$ に替え、
三生成元のそれぞれに $e$ の倍数を加える。
平行移動の係数は、まず $XY-qYX$ の $X,Y$ 係数、
次に $ZX-qXZ$ の $X$ 係数を消すことで一意に決まる。
最後の移動は $YZ-qZY$ の $Y$ 係数も同時に消す。
こうして得た $x,y,z$ は節点の三関係を満たし、
\[
 P_x=-\lambda(q^3-1)/q,\qquad
 P_y=-(q^3-1)/\lambda^2,\qquad
 A_z=-\lambda(q^2-1)/(q\Delta)
\]
である。さらに命題\ref{prop:cubic-CY-three-alpha-nodal-secant-center}
の $t=1$ での $J$ は $e$ のスカラー倍となる。
$D=-J/(q^3-1)$ のスカラー値を切断係数とする。

係数計算では $b^{-1}=b^2-f_1b+f_2$、
$c^{-1}=c^2-g_1c+g_2$ を使う。
上文献第8.1節の四書換えの先頭語は $bbb,ccc,bcbc,bbcc$ である。
例えば $bcbc=e_2bc-e_1+q c^{-1}b^{-1}$、
$cbcb=e_2cb-e_1+q b^{-1}c^{-1}$ から後二つを得る。
この書換えのもとで三関係と中央三次のスカラー恒等式を
七変数の有理関数体上で直接検算する。
最終式の分母を検査することで全単純根チャートへ特殊化する。
全係数と再現手順は
\path{verify_cubic_CY_three_alpha_nodal_E6.py} と対応するJSONにある。

次に上の準同型が同型であることを示す。
同論文定理6.5の $F_{123}$ 濾過では、$c$ は次数零、$b$ は次数1である。
球面部分の随伴環は、節点有理曲線の次数1の直線束と
$q$ 作用による捩れ斉次座標環である。
規格化座標を $\eta$ とすると、次数 $n$ は
\[
 V_n=\bigl\{F(\eta):\operatorname{div}_\infty F
        \le \textstyle\sum_{i=0}^{n-1}[q^i],\ F(0)=F(\infty)\bigr\}.
\]
その積は次数 $m$ の右因子に対して左因子を $\eta\mapsto q^{-m}\eta$ で引き戻す。
$c$ の伴随行列を $C$ とすれば、同定の下で $b$ の行列は
\[
 \frac{\eta E_{11}C-CE_{11}}{\eta-1}
 =\frac{\eta E_{13}-E_{21}}{\eta-1}.
\]
$e$ は $C$ の固有直線への射影である。
$C$ の右固有列と左固有行を $v,w$ とすると
$v_3=\lambda v_1$、$w_2=\lambda w_1$ であり、
$v_1w_1\ne0$、$wv\ne0$ である。
従って $eE_{13}e=eE_{21}e\ne0$ であり、$eE_{23}e\ne0$ でもある。
このため $X$ の次数1の像は非零定数、$Y$ の次数2の像は
\[
 \frac{\eta}{(q^{-1}\eta-1)(\eta-1)}
\]
の非零倍となる。平行移動はこの結論を変えない。

これら二つが $\bigoplus_n V_n$ を生成する。
次数 $n\ge3$ で $q^{n-1}\ne1$ なら、次数1の定数を
左右から $V_{n-1}$ に掛けた二つの像は、
異なる一箇所ずつの極を除く余次元1部分空間なので和が $V_n$ である。
$q^{n-1}=1$ なら $n\ge5$ であり、両者が取り逃すのは $\eta=1$ の最高極次数である。
次数2の上の関数を $V_{n-2}$ の右から掛ければ、同点で最高極次数を持つ元を得る。
実際、右からの積で左因子の極は $q^2,\ldots,q^{n-1}$ へ移り、
その中の $\eta=1$ に最大次数の極を持つ元を選べる。
次数 $n-2\ge3$ なので、極を一つ減らす条件は真に余次元1である。
これで全位数について帰納が閉じる。
従って $eHe$ は $x,y$、従って $x,y,z$ で生成される。

$x,y,z$ に重み $(1,2,3)$ を与えた節点アフィン三関係の環は
PBW基底 $x^iy^jz^k$ を持つ。書換えで各項の重みは増えない。
随伴環をさらに重み $(5,6,8)$ で濾過すると三変数量子多項式環となるから、
最初の随伴環も整域である。
中央三次の先頭重みは6であり、$xyz$ の係数が非零なので正則。
従って中央切断の重み付きHilbert級数は
$(1-t^6)/((1-t)(1-t^2)(1-t^3))=1+t/(1-t)^2$ となる。
生成元の $F_{123}$ 次数は高々 $1,2,3$。
上の随伴環の生成証明は準同型が濾過に関して全射であることも示す。
同じHilbert級数なので同型である。

最後に $(f_1,f_2,g_1,\lambda,e_1,e_2)=(0,0,0,1,0,0)$ で
六係数変化と三つの生成元倍率変化を取る。
係数順を $(P_x,P_y,A_z,A_x,A_y,B_x,B_y,B_z,D)$ とするとヤコビアンは
\[
 \frac{(q-1)^{17}(q+1)^6(q^2+q+1)^3}{59049q^8}.
\]
この恒等式も同じ検証器で再現する。全許容 $q$ で非零なので支配性が従い、
係数像は稠密開集合を含む。
\end{proof}

\begin{corollary}\label{cor:cubic-CY-three-alpha-nodal-E6-H3}
節点型には、各固定有限位数 $q\ne0,\pm1$、$q^3\ne1$ で
非空な稠密開集合のH3切断が存在する。
また各固定非捩れ $q$ で、前定理のヤコビアンを評価した点はH3である。
\end{corollary}
\begin{proof}
有限位数の場合、一般パラメータの $H$ は滑らかな中心上のAzumaya環である
\cite[定理5.6、5.9]{EtingofOblomkovRainsGDaha}。
その中心は既約である \cite[命題6.6]{EtingofOblomkovRainsGDaha}。
非零冪等元 $e$ の階数は局所一定なので $e$ はfullで、
$eHe$ も同じ中心上Azumayaで大域次元2となる。
この一般パラメータ域と定理の単純根域は交わり、支配写像によりその像も稠密開集合を含む。
実際、各スペクトル三つ組の積を固定した族でも、
定理5.4の行列式条件で恒等的に定数となる積は各三つ組の重複度が一定の場合に限られ、
その場合の表現次元 $p<3N$ は $q^{p/3}=1$ を満たせない。

非捩れの場合、指定点の三つの定義多項式は純三乗を定める。
Etingof--Oblomkov--Rainsの第2節
\cite{EtingofOblomkovRainsGDaha}により $H$ は量子トーラスと
位数3の回転の半直積である。
非捩れ量子トーラスは単純で、回転の非自明な冪は外部自己同型である。
実際単元は尺度倍のLaurent単項式に限られ、
内部自己同型は格子指数を変えない。
単純環の外部有限群作用の半直積は単純なので $e$ はfull。
量子トーラスは一変数Laurent環からのskew Laurent拡大で大域次元は高々2、
有限群平均化により半直積と $eHe$ も大域次元は高々2である。
両場合とも、節点型の整域・rank-oneと射影境界の既出判定
\ref{thm:cubic-CY-three-alpha-quantum-boundary-all-q}からH3を得る。
\end{proof}

\begin{proposition}\label{prop:cubic-CY-three-alpha-nodal-E6-finite}
定理\ref{thm:cubic-CY-three-alpha-nodal-E6}の表示を持つ有限位数 $N$ の切断では、
アフィン環 $eHe$ は有限生成可換中心環上有限である。
H3の必要十分条件は、その次元 $3N$ 以下の全単純対で第三Extが消えることである。
\end{proposition}
\begin{proof}
$F_{123}$ 随伴環の中心に、次数 $mN$ のスカラー行列で
$\eta\mapsto q\eta$ に不変なものからなる部分環を取る。
積の引き戻しはこの部分環の因子には現れない。
これは $r=\eta^N$ とした節点有理曲線の
次数1直線束の斉次座標環である。
実際、極は $r=1$ に高々 $m$ 位、両節点枝での値は等しい。

規格化座標の $r=\eta^N$ は二つの節点曲線の間の有限射を与え、
$\sum_{i=0}^{N-1}[q^i]$ は $[1]$ の引き戻しである。
従って各次数剰余 $j\bmod N$ の
$\bigoplus_{m\ge0}V_{j+mN}$ は
有限射による直線束の押し出しの切断加群であり、
上の中心部分環上有限となる。
Etingof--Oblomkov--Rainsの定理6.7
\cite{EtingofOblomkovRainsGDaha}により随伴中心は元の中心から持ち上がる。
有限個の中心生成元と有限個の加群生成元を持ち上げれば、
次数帰納で元の球面環の中心有限性が従う。

同論文定理5.3の次数 $3N$ のPI恒等式は部分環 $eHe$ も満たす。
従って単純の次元は高々 $3N$ である。
系\ref{cor:cubic-CY-three-alpha-triangle-finite-effective}の
有限可換環上の完備化と単純対Extの議論に、
節点の射影境界判定を組み合わせて同値を得る。
\end{proof}

この定理は単純根チャートから生じる各アフィン環を同定する。
以下で全許容 $q$ の有限大域次元と、係数写像の全射性を順に決める。

\begin{theorem}\label{thm:cubic-CY-three-alpha-nodal-E6-separable}
定理\ref{thm:cubic-CY-three-alpha-nodal-E6}のチャートで三つの
定義三次多項式がすべて平方自由であるとする。
有限位数・非捩れの双方で、対応する切断がH3であるための必要十分条件は
$H/HeH=0$ である。
この商は17個以下の語で張られ、その消滅は有限の階数計算で判定できる。
各許容 $q$ で節点切断のH3域は稠密開集合を含む。
\end{theorem}
\begin{proof}
$\kappa^3=q$ とし、$a=\kappa(bc)^{-1}$ を導入する。
$H$ は三つの三次可換代数 $A_a,A_b,A_c$ の自由積に
$abc=\kappa$ を課した環である。
三つの平方自由性により $A_i$ は分離代数である。

まず $H$ は大域次元高々2である。分離代数の非可換微分を
$\Omega^1 A_i=\ker(A_i\otimes_k A_i\to A_i)$ とする。
自由積の微分は各 $\Omega^1 A_i$ を誘導した直和である。
商の微分を取ると、次の双加群列が左端以外で完全になる。
\[
 H\otimes_k H
 \xrightarrow{d_2} \bigoplus_{i=a,b,c}
       H\otimes_{A_i}\Omega^1 A_i\otimes_{A_i}H
 \xrightarrow{d_1} H\otimes_k H \longrightarrow H\longrightarrow0 .
\]
実際、誘導微分の直和から $H$ の微分への商は
$d(abc-\kappa)=(da)bc+a(db)c+ab(dc)$ で生成される。
これは自由積上の任意の導分が、$abc-\kappa$ の導分を消せば
$H$ へ一意に降りることからも従う。
$d_2$ の三成分はこの三項の左右倍であり、
$d_1$ は通常の微分の包含である。
$A_i$ が分離的なので各成分は $H\otimes_k H$ の双加群直和因子である。

$T_g$ を $u\otimes v\mapsto ug\otimes g^{-1}v$ と書く。
各成分を $H\otimes_k H$ に埋め込んで見ると、
$d_2/\kappa$ の三成分は $1-T_a,\ T_a-T_{ab},\ T_{ab}-1$ である。
その共通核の元の第一テンソル因子の係数空間は、
$a,ab$ の右乗法で不変である。
従って $a,b,c$ の右乗法で不変な有限次元右イデアルになる。
一方、定理6.5の $F_{123}$ 随伴環には次数1のスカラー恒等行列があり、
その右乗法は単射である。
その持ち上げの冪を右から掛けると、非零右イデアルの一元から
任意に大きい相異なる濾過次数の元が生じる。
従って有限次元右イデアルは零で、$d_2$ は単射となる。
これで長さ2の双加群射影解消を得て、任意の左加群の射影次元は高々2である。

次に $Q_0=H/HeH$ を考える。
$e=0$ を課すことは $c$ の定義多項式を
$c^2+(\lambda-g_1)c+\lambda^{-1}$ に置き換えることに等しい。
この二次式と $b$ の三次式から $cc,bbb$ を消去できる。
また三次式 $(bc)^3-e_2(bc)^2+e_1bc-q=0$ により $bcbc$ と $cbcb$ を
より短い語に直せる。
$b^{-1}c^{-1}$ と $c^{-1}b^{-1}$ の二乗を同じように展開すれば、
$b^{-1}$ の最高語が $bb$、$c^{-1}$ の最高語が $c$ なので
$bbcbbc,cbbcbb$ も短い語に直せる。
これら六語を部分語に持たない語は $c$ を高々二つしか含まない。
$c$ が零個、一個、二個の場合の数はそれぞれ $3,9,5$ である。
よって $\dim Q_0\le17$ である。

この語集合を基底にした空間に、
指定順の書換えで $b,c$ の左乗法行列を作る。
三つの定義関係をこれらの行列に代入して得る像を取り、
$b,c$ で不変な最小部分空間まで閉じる。
その商は $Q_0$ の巡回加群であり、各基底語は $1$ から得られる。
一方、書換えは元の関係で成立するので、この商から $Q_0$ への全射がある。
巡回ベクトルからの射と互いに逆であり、これで商の次元を判定できる。
この有限商計算は \path{verify_cubic_CY_three_alpha_nodal_E6_full.py} で再現する。

$Q_0=0$ なら $e$ はfullなので $eHe$ は大域次元高々2。
$Q_0\ne0$ なら、非零左 $Q_0$ 加群 $M$ と $M$ への非零な
$H$ 線型写像 $g:Q_0\to M$ を取る。
上の双加群解消を $M$ でテンソルして $Q_0$ へのHomを取ると、
最高項は $\operatorname{Hom}_k(M,Q_0)$ である。
そこへの微分の像は、各成分が双加群直和因子であることと
$d_2$ の式から、同時共役による差の和である。
汎関数 $f\mapsto\operatorname{tr}_{Q_0}(f\circ g)$ は
この差をすべて消す。実際 $g$ が $H$ 線型なので、
同時共役は $f\circ g$ の通常の共役になる。
$g\ne0$ なので、この汎関数は全空間上では非零である。
従って $\operatorname{Ext}^2_H(M,Q_0)\ne0$ となる。
補題\ref{lem:cubic-CY-three-alpha-defective-idempotent}により
$eHe$ の大域次元は無限である。
節点の射影境界判定から、この条件はH3判定に一致する。

最後に $Q_0=0$ は上の有限行列の階数の開条件である。
一般スペクトルで $Q_0\ne0$ とすると、次元17以下の単純から
三つのスペクトル根の重複度による行列式の等式が得られる。
各三つ組の積を固定しても、消した $c$ の根の重複度が零なので、
正の表現次元で恒等式となる組はない。
有限個の等式を避ければ商は零である。
三多項式の分離域は稠密であり、定理
\ref{thm:cubic-CY-three-alpha-nodal-E6}の支配性により
このH3域の係数像は各 $q$ で稠密開集合を含む。
\end{proof}

\begin{proposition}\label{prop:cubic-CY-three-alpha-nodal-E6-full-repeat}
定理\ref{thm:cubic-CY-three-alpha-nodal-E6}のチャートで $H/HeH=0$ とする。
切断がH3であるための必要十分条件は、三つの定義多項式が相異なる根を持つことである。
従って単純根表示内で、上の定理と合わせても未判定となる非捩れ係数は、
少なくとも一つの多項式が重根を持ち、かつ $H/HeH\ne0$ の場合に限られる。
\end{proposition}
\begin{proof}
三つの三次部分代数を $A_i\subset H$ とする。
まず $A_c=k[c]$ 上では、$F_{123}$ の各正次数商は
$\operatorname{Mat}_3(k)$ の直和である。
実際、定理6.5の行列値関数において、許容極の主部と無限遠の値を取ると、
$F(0)=CF(\infty)C^{-1}$ は一つの行列係数から解ける。
$C$ の多項式による左右乗法はこの関係を保つので、この同定は
$A_c$ 双加群として成り立つ。
$C$ は巡回伴随行列で、その作用空間は正則 $A_c$ 加群である。
有限単項式商 $A_c$ はFrobenius環なので、その双対も正則加群に同型。
従って $\operatorname{Mat}_3(k)$ は $A_c\otimes A_c^{\rm op}$ 上自由である。

濾過の各正次数商はこの双加群圏で射影的なので、各段階の完全列は分裂する。
従って $H$ は左右それぞれ $A_c$ 上自由で、かつ $A_c$ を双加群直和因子に持つ。
$abc=\kappa$ は巡回置換で保たれるため、
$A_a,A_b$ についても同じ結論が成り立つ。
各 $A_i$ 加群 $M$ に対し、$H\otimes_{A_i}M$ を $A_i$ に制限すると
$M$ が直和因子になる。また $H$ 射影加群は制限しても射影的である。
従って $H$ の大域次元が有限なら、各 $A_i$ の大域次元も有限になる。
可換三次代数 $A_i$ が重根を持てば非半単純Artin環であり、
その被約単純の射影次元は無限である。
よって重根があれば $H$ の大域次元は無限。
$e$ がfullなので $eHe$ も同じである。
平方自由の場合は定理\ref{thm:cubic-CY-three-alpha-nodal-E6-separable}による。
\end{proof}

\begin{theorem}\label{thm:cubic-CY-three-alpha-nodal-E6-complete-corner}
定理\ref{thm:cubic-CY-three-alpha-nodal-E6}の全単純根チャートにおいて、
切断がH3であるための必要十分条件は、三つの定義三次多項式が平方自由で、
かつ17語で判定する商 $H/HeH$ が零であることである。
有限位数・非捩れの双方に適用できる。
\end{theorem}
\begin{proof}
既出の判定から、残るのは三次多項式に重根があり
$Q_0=H/HeH\ne0$ の場合の排除である。
ここでは相対双加群列と有限次元加群の双対を用いる。

最初に $H$ は左noetherianであり、$eH$ は $B=eHe$ 上有限生成である。
$F_{123}$ 随伴環において、前者の行列値切断と後者の左から $e$ を掛けた切断は、
いずれも規格化上で自明なベクトル束の固定貼合せから得られる。
次数 $n$ の極の因子を $D_n=\sum_{i=0}^{n-1}[q^i]$ とする。
スカラー行列の次数2の部分は、定数 $1$ と
$\eta/((\eta-1)(\eta-q))$ による基点のないpencilを持つ。
その引き戻しを使い、次数 $n$ の切断に左から掛ける写像は
$n\ge3$ で次数 $n+2$ へ全射である。
実際、pencilの完全列をベクトル束と $D_n$ で捩ると、
核の直線束部分の次数は $n-2$ である。
規格化の $0,\infty$ で値を独立に取れるので、
貼合せ完全列によりこの核の第一コホモロジーは零となる。
従って行列値随伴環とその左からの $e$ の像はスカラー座標環上有限生成。
スカラー座標環は次数1の直線束による節点曲線の捩れ斉次座標環であり、
曲線上の振幅条件によりnoetherianである
\cite{AVdBTwistedHomogeneous}。
従って随伴環と $H$ は左noetherianとなり、
球面随伴環の生成元を持ち上げると $eH$ は左 $B$ 上有限生成となる。

三つの部分代数 $A_i$ は、根の分離を仮定しなくてもFrobenius環である。
命題\ref{prop:cubic-CY-three-alpha-nodal-E6-full-repeat}により
$H$ は各 $A_i$ 上左右自由である。
従って定理\ref{thm:cubic-CY-three-alpha-nodal-E6-separable}の
相対双加群列の各微分成分は $H\otimes_k H$ に単射で入る。
有限次元右イデアルを排除した同じ証明で $d_2$ は単射である。
中間完全性は微分の普遍性から従うので、この相対列は重根がある場合も完全である。
さらに $\Omega^1 A_i$ は片側 $A_i$ 加群として自由加群の直和因子である。
従ってこの列は右 $H$ 加群として項ごとに射影的で、終項 $H$ も射影的。
任意の左 $H$ 加群 $M$ でテンソルしても完全性を保つ。

$M$ を有限次元とする。右辺の係数を $H$ としたとき、
この相対列の各項は高次Extを持たない。
自由誘導項では明らかであり、各 $A_i$ からの誘導項では
$H$ が $A_i$ 上自由従って入射的だからである。
ここで有限次元Frobenius環上では自由加群は入射的であり、
無限直和の場合もBaerの判定で同じ結論を得る。
従って $j>2$ で $\operatorname{Ext}^j_H(M,H)=0$ である。
また最高微分への像は
\[
 f\longmapsto f-g f(g^{-1}{-}),\qquad g=a,ab
\]
の像の和である。
実際、各微分成分から $H$ への写像は、
$A_i$ 上の入射性により誘導前の零化微分加群の包含を通して
$A_i\otimes_k M$ 全体への写像に延長できる。
よって $d_2$ の三つの差の式をそのまま使える。
$\operatorname{Hom}_k(M,H)=M^*\otimes_k H$ で、
これらの差は $vg\otimes h=v\otimes gh$ を課す。
$a,ab$ と逆元は $a,b,c$ を生成するから
\[
 \operatorname{Ext}^2_H(M,H)\simeq M^*\otimes_H H=M^*
\]
となる。これは右 $H$ 加群としての同型である。

$B$ の大域次元が有限と仮定すると、
$eH$ は有限生成射影 $B$ 加群による有界解消を持つ。
従って三角
\[
 He\otimes_B^{\mathbf L}eH\longrightarrow H\longrightarrow C
 \longrightarrow (He\otimes_B^{\mathbf L}eH)[1]
\]
の $C$ は左 $H$ 上完全複体となる。
$C$ のコホモロジーは有限生成で、$e$ で零になる。
$Q_0$ は有限次元なので、これらは有限次元加群である。
$H^0(C)=Q_0\ne0$、正次数は零であり、最低の非零次数を $-n$ とする。
上の高次Ext消滅から、標準切断により
$C^\vee=\mathbf R\operatorname{Hom}_H(C,H)$ の最高コホモロジーは
次数 $n+2$ の $(H^{-n}(C))^*$ である。
$C$ は完全なので $\mathbf R\operatorname{Hom}_H(C,C)\simeq
C^\vee\otimes_H^{\mathbf L}C$。
その次数 $n+2$ のコホモロジーは
$(H^{-n}(C))^*\otimes_H Q_0$ であり、
双対も $Q_0$ 加群なので非零である。
一方、$eC=0$ と三角の導来随伴から
$\mathbf R\operatorname{Hom}_H(C,C)\simeq C$。
正次数の非零は矛盾する。従って $B$ は無限大域次元でありH3ではない。
\end{proof}

\begin{theorem}[節点型の全係数の球面認識とH3判定]
\label{thm:cubic-CY-three-alpha-nodal-E6-surjective}
$q\ne0,\pm1$、$q^3\ne1$ とする。
$P_xP_yA_z\ne0$ の全節点型中央三次切断は、定理
\ref{thm:cubic-CY-three-alpha-nodal-E6}の単純根チャートから得られる。
任意のそのような表示で、三つのスペクトルが平方自由で
$H/HeH=0$ であることがH3の必要十分条件である。
\end{theorem}
\begin{proof}
定理\ref{thm:cubic-CY-three-alpha-nodal-E6}の生成元 $(x,y,z)$ を
$(\lambda x/A_z,y/A_z,\lambda z/A_z)$ に替える。
すると $(P_x,P_y,A_z)=(-(q^3-1)/q,-(q^3-1),1)$ となる。
対角倍率で非零の任意の三係数へ移れるので、この規格化で証明すればよい。
残る六係数は次の二変数 $U,V$ を通じて $g_1,\lambda$ に依存する。
\[
 U=g_1+(q-1)\lambda,\qquad V=q\lambda(g_1-\lambda)+\lambda^{-1}.
\]
実際、$c$ の旧スペクトルの選択根を $\lambda$ から $q\lambda$ に替えた
三つ組は $T^3-UT^2+VT-q$ の根である。
このため $g_1=U+(1-q)\lambda$ と置き、
$q^2\lambda^3-qU\lambda^2+V\lambda-1=0$ で六係数を約分すると、
$f_1,f_2,U,V,e_1,e_2$ の多項式になる。

$\theta^3=q$ を一つ選ぶ。
三組の積が1の固有値を持つトーラスで、各組の自然表現と双対表現の
指標を順に
\[
 (f_1,f_2),\qquad (U/\theta,V/\theta^2),\qquad
 (e_1/\theta^2,e_2/\theta)
\]
とする。三組目は規格化された $bc$ の逆元の三つ組に対応する。
各 $A_2$ の自然表現の重みを $w_1,w_2,w_3$ と書き、
内積を $(w_i,w_i)=2/3$、$(w_i,w_j)=-1/3$ ($i\ne j$) とする。
$E_6$ の単純根の順序を、三組を貫く $(-w_1,-w_1,-w_1)$、
第1組の $w_1-w_2,w_2-w_3$、第2組の同じ二根、
第3組の $w_1-w_2$ とする。そのCartan行列は $E_6$ 型である。
三組の重み格子の直和内で、これらの六根と整数の内積を持つ
部分格子 $\mathcal P$ は指数3であり、$E_6$ の重み格子に一致する。
実際、六根のCartan行列式は3であり、
三つの $A_2$ 重み格子の直和の共体積は $3^{-3/2}$ である。
整数内積で定まる指数3の部分格子と $E_6$ 重み格子の共体積は
ともに $3^{-1/2}$ となる。

基本重みを $\omega_i$、支配重み $\nu$ の軌道和を $m_\nu$ と書く。
六係数のLaurent展開は六つの単純反射に不変である。
下表は各係数の、分母を除いた軌道和展開の支持である。
$d=(q-1)(q+1)^2$ と置いた。
\[
\begin{array}{c|c|c|c|l}
\text{係数}&\text{分母}&\text{基本重み}&\text{分子係数}
 &\text{ほかの支配重み}\\ \hline
A_x&1&\omega_3&\theta^2&\\
A_y&1&\omega_5&\theta^4&\\
B_x&d&\omega_2&-\theta^7&\omega_5\\
B_y&d&\omega_4&-\theta^8&\omega_3\\
B_z&d&\omega_6&-\theta^6&0\\
D&q(q-1)^3(q+1)^3&\omega_1&\theta^{15}&0,\omega_6,\omega_3+\omega_5
\end{array}
\]
さらに $m_{\omega_3}m_{\omega_5}=m_{\omega_3+\omega_5}+6m_{\omega_6}+27$。
これらの恒等式は、三組の指標をそれぞれ三つの重みの和に展開し、
各単純反射で係数を比較すれば得られる。
\path{verify_cubic_CY_three_alpha_nodal_E6_surjective.py}
は、$\lambda$ の消去と軌道和の係数を有理数上で再現する。

基本重みの軌道和は $k[\mathcal P]^{W(E_6)}$ を自由生成する。
実際、支配重み $\nu=\sum_i a_i\omega_i$ に対し
$\prod_i m_{\omega_i}^{a_i}$ の最高項は $m_\nu$ で係数1、
残余は支配順序で真に下にあるため帰納法で従う。
表の主要係数は全許容 $q$ で単元なので、規格化した六係数も
同じ不変式環を自由生成する。
従ってトーラスからこの六係数空間への射は、
指数3の格子包含による有限全射と有限Weyl群による商の合成であり全射である。

最後に、この全射の各持上げに対応する $T^3-UT^2+VT-q$ の根の多重集合から
$\mu$ を一つ選び、$\lambda=\mu/q$ と戻す。
$\lambda$ が旧 $c$ の単純根となる選択が存在する。
存在しなければ、全ての異なる根 $\mu$ に対して $\mu/q$ も
同じ根集合に入る。非零の高々三点が $q^{-1}$ 倍で巡回することになり、
$q,q^2,q^3\ne1$ に反する。
こうして $\Delta\ne0$ の表示が得られ、球面同型と
定理\ref{thm:cubic-CY-three-alpha-nodal-E6-complete-corner}の判定を全係数に適用できる。
\end{proof}

\begin{theorem}[節点型の許容域は $E_6$ 判別式の補集合]
\label{thm:cubic-CY-three-alpha-nodal-E6-root-discriminant}
前定理の規格化係数空間 $T_{E_6}/W(E_6)$ において、
H3域は根指標 $t^\alpha=1$ の超曲面の像を除いた部分である。
この判別式は既約である。
\end{theorem}
\begin{proof}
三つの積1スペクトルを $(b_i),(c_i),(d_i)$ とする。
前定理の根は、各組内の比に当たる18個と、
$b_i c_j d_k$ またはその逆数に当たる54個である。
これは三重みの和と組内の重みの差の72個を
一つの単純根から単純反射で生成する有限計算でも確認できる。
従って根超曲面上の点はWeyl群で $b_1=b_2$ を満たす点に移る。
その点は同一の六係数を表し、$b$ に重根がある。
前定理の全単純根チャートでの必要十分判定によりH3ではない。
全ての根が一軌道をなし、単純根は重み格子の原始元なので、
一つの既約トーラス超曲面の像が判別式全体となる。

逆に全ての根指標が1と異なるとし、第二組の一つを旧スペクトルに戻して
$c$ の単純根 $\lambda$ を得る。
旧三スペクトルは全て平方自由である。
$H/HeH\ne0$ なら有限次元であるため単純加群 $V$ が存在する。
$H$ の生成元の作用を使って、$V$ 上の三行列
$b,\theta^{-1}c,\theta(bc)^{-1}$ を取る。その積は1であり、
固有値は順に $(b_i)$、選択した一つを除く $(c_i)$、$(d_i)$ に属す。
いずれも半単純で、三行列は既約な組をなす。

$n=\dim V$、固有値重複度を $m_{\ell j}$ とする。
三共役類の積写像の微分の像の跡直交補は
三行列の共通中心化環であり、既約性からスカラーに等しい。
従って微分の階数は $n^2-1$。
積が1のファイバーの接空間は同時共役軌道の接空間を含むので、
\[
 \sum_{\ell=1}^3\left(n^2-\sum_j m_{\ell j}^2\right)
 \ \ge\ 2(n^2-1).
\]
第一、第三組は三つ、第二組は二つの固有値を許すため、
$\sum_{\ell,j}m_{\ell j}^2\ge 7n^2/6$ である。従って $n\le3$。
$n=1$ では行列式の積が $b_i c_j d_k=1$ を与える。
$n=2$ では上の不等式から三組とも二固有値を一回ずつ使う。
各組の全固有値の積は1なので、補三つ組が $b_i c_j d_k=1$ を与える。
$n=3$ では第一、第三組は三固有値を一回ずつ使い、
第二組は残した二固有値を2回と1回使う。
第二組の行列式が1なので、二回用いた固有値は
選択して除いた固有値に等しい。どの場合も根指標の仮定に反する。
従って $H/HeH=0$ でありH3である。
\end{proof}

\begin{corollary}\label{cor:cubic-CY-three-alpha-nodal-E6-open}
節点型のH3域は規格化係数空間においてZariski開であり、
各許容 $q$ のファイバーで稠密である。
全幾何学的矢で割った節点型の層は既約で次元7、一般慣性は自明である。
\end{corollary}
\begin{proof}
判別式はトーラスの真閉集合の有限商像である。
$q$ も動かす場合は $\theta^3=q$ の有限全射を最初に加えれば、
同じ構成から閉性が従う。
規格化後の全幾何学的矢は次命題の有限群作用なので、
次元と既約性を保ち、一般慣性も次命題から従う。
\end{proof}

\begin{proposition}\label{prop:cubic-CY-three-alpha-nodal-arrows}
節点型の $P_xP_yA_z\ne0$ は $P_x=P_y=A_z=1$ に規格化できる。
切断を $J/(q^3-1)+Dt^3$ と書く。
共通スカラーを除いた全幾何学的矢は、次の対角変換と交換変換で生成される。
\[
\begin{array}{c|c|c}
 & (x,y,z)\text{ の像} & q'\\ \hline
 \text{対角} & (\zeta x',\zeta^2y',z'),\quad \zeta^3=1 & q\\
 \text{交換} & (q^{-2}y',q^{-2}x',-q^{-3}z') & q^{-1}
\end{array}
\]
ここで $t\mapsto t'$ とできる。対角変換の係数作用は
\[
 (A_x,A_y,B_x,B_y,B_z,D)
 \longmapsto(\zeta^2 A_x,\zeta A_y,\zeta B_x,\zeta^2 B_y,B_z,D)
\]
であり、交換変換では
\[
 (A_x,A_y,B_x,B_y,B_z,D)
 \longmapsto
 \left(q^2 A_y,q^2 A_x,q^4 B_y,q^4 B_x,-q^3 B_z,
       -q^6\left(D-\frac{B_z}{q^3-1}\right)\right).
\]
許容切断へ制限すれば全点付き群oidとなる。
共通スカラーを除いた慣性は、$A_x=A_y=B_x=B_y=0$ で $\mu_3$、
その他では自明である。
\end{proposition}
\begin{proof}
$(x,y,z)$ の対角倍率による $P_x,P_y,A_z$ の三つの変換式は独立で、
代数閉体上で三つを1にできる。
固定出典の節点同型式(N1)--(N2) \cite{AS4CYpinned}により、
全生成元変換は平行移動を持たない対角変換か接線交換であり、
共通スカラーで $t\mapsto t'$ とできる。
純平方の係数と $xy-qyx$ の $z$ 係数を比較すると上の倍率だけが残る。
他の係数は三関係への代入で得られる。

交換で $J/(q^3-1)$ の順序語 $xyz$ の係数は $-q^{-6}$ 倍になる。
$-(q^3-1)^{-1}txy$ の交換を順序化すると
$-(q^3-1)^{-1}B_z t'^3$ が生じる。
他に純 $t'^3$ 項は生じないので、切断の同じ直線を送る条件が
$D$ の上記変換である。
この全係数恒等式は \path{verify_cubic_CY_three_alpha_triangle_center.py} で再現する。
それぞれの表示は逆に包絡と切断線を移す。
反Serre環から包絡と切断線を復元できるので点付き矢は尽くされる。
$q\ne q^{-1}$ から交換は一つの点を固定せず、
三乗根の対角作用の固定条件から慣性が従う。
\end{proof}

\begin{corollary}\label{cor:cubic-CY-three-alpha-nodal-E6-relative-coordinates}
$Q_{\mathrm{nod}}=\{q\in\mathbb G_m:q^2\ne1,\ q^3\ne1\}$ とし、
$T_{\mathrm{reg}}\subset T_{E_6}$ を全根指標が1と異なる開集合とする。
全許容の規格化係数空間は $Q_{\mathrm{nod}}\times (T_{\mathrm{reg}}/W(E_6))$
と代数的に同型である。
ここで三乗根の対角変換は $T_{E_6}$ の中心 $\mu_3$ の乗法、
接線交換は $(q,t)\mapsto(q^{-1},t^{-1})$ に対応する。
従ってこの係数表示と矢の商は
\[
 \left[
   \bigl(Q_{\mathrm{nod}}\times (T_{\mathrm{reg}}/W(E_6))\bigr)
   /(\mu_3\rtimes C_2)
 \right]
\]
という代数的表示スタックを持つ。
\end{corollary}
\begin{proof}
定理\ref{thm:cubic-CY-three-alpha-nodal-E6-surjective}の規格化から
$P_x=P_y=A_z=1$ へ移す倍率は、
$C=q^3-1$ として $(-\theta/C,-\theta^2/C,q/C^2)$ である。
この倍率を軌道和展開に適用すると、各 $\theta$ 指数は3の倍数となる。
例えば $m_{\omega_i}$ を $m_i$ と書けば、
\[
\begin{aligned}
 A_x&=q^2m_3/C^2,\qquad A_y=q^2m_5/C^2,\\
 B_x&=\frac{q^4m_2+(-q^5+6q^4-q^3)m_5}{C^3(q-1)(q+1)^2},\\
 B_y&=\frac{q^4m_4+(-q^5+6q^4-q^3)m_3}{C^3(q-1)(q+1)^2}.
\end{aligned}
\]
全六式とその三角的な逆写像は $Q_{\mathrm{nod}}$ 上正則である。
$\omega_3+\omega_5$ の軌道和は前定理の積恒等式で消去できる。
従って根判別式を除いた後も全 $q$ で積の表示を得る。

三組の固有値における中心元の代表を
$(1,\zeta,\zeta^{-1})$ と取ると、その全根指標は1で、
六係数に順に $(\zeta^2,\zeta,\zeta,\zeta^2,1,1)$ が掛かる。
反転 $t\mapsto t^{-1}$ は基本軌道和の順序を
$m_2\leftrightarrow m_4$、$m_3\leftrightarrow m_5$ と交換し、
$m_1,m_6$ を固定する。
$q\mapsto q^{-1}$ と合わせた係数変換を六式で計算すると、
前命題の接線交換と一致する。特に切断定数の $B_z$ 補正も一致する。
\path{verify_cubic_CY_three_alpha_nodal_E6_surjective.py}
は指数条件、中心の作用と反転後の全六式を検算する。
両作用の間の関係は中心の逆元作用なので、半直積による商を得る。
一般の内在的Abel圏の族の相対再構成には、さらに相対認識定理を要する。
\end{proof}
